Vision-language-action (VLA) policies have shown strong promise as generalist
controllers for robot manipulation, yet their action predictions can remain
fragile under changes in visual perspective. We identify a systematic grounding
failure: lateral camera shifts induce co-directional shifts in predicted
end-effector trajectories, even when the underlying robot-scene state is
unchanged. This behavior suggests that current VLAs may exploit image-space
shortcuts, such as pixel-coordinate correlations, rather than grounding actions
in proprioceptive state and robot-object geometry.
We introduce \textsc{ProprioVLA}, a self-supervised framework for
proprioceptive grounding in VLA policies. Our key insight is that the
end-effector provides a natural robot-kinematic anchor between visual evidence
and proprioceptive state. A grounded policy should not merely recognize task
objects in the scene, but should also understand where its own manipulator is
and how its visual motion relates to the robot's underlying kinematics. To
instill this grounding bias, \textsc{ProprioVLA} introduces two complementary
objectives during policy fine-tuning: a visual grounding objective that
encourages the policy to attend to the end-effector, and a proprioceptive
reconstruction objective that requires the grounded representation to recover
the end-effector state.
Crucially, \textsc{ProprioVLA} obtains grounding supervision from robot
interaction itself. Rather than relying on static gripper appearance,
\textsc{ProprioVLA} exploits a motion-kinematic consistency principle: the true
end-effector region is the one whose visual motion is both locally salient and
consistent with proprioceptive change. This enables the automatic mining of
confidence-weighted end-effector targets with uncertainty-aware filtering,
without manual annotations, external grounding datasets, or supervision from
stronger foundation models. Experiments on manipulation benchmarks and
controlled camera-perturbation evaluations show that \textsc{ProprioVLA}
improves task success and robustness over standard VLA fine-tuning, suggesting
that proprioceptive self-grounding is a practical route toward spatially
reliable VLA control.